购买系统跨地域运行时，业务身份沿着一条完整的事实链移动：入口找到购买所属分区，分区保存受理，库存分区形成联合预留决定，工作者取得付款资格，平台保存收款，购买分区固定原结果与事件，消费者接受履约，发货平台产生实际效果。地域部署决定这条链上哪些通信需要跨域、哪些保存域可以在故障后继续，以及恢复者从哪里找回原来的责任。

沿用前五章的三个逻辑分区 b1、b2、b3，当前分别由 G1、G2、G3 承载。每组是完整的 Raft 状态机，包含会话、$S$ 绑定、原结果、C/T、Outbox、退款占额以及本组承担的库存或协调事实。控制组单独保存归属、迁移和配置资格。图~\ref{fig:ch21-regions} 中，每一组在 A、B、C 三个地域各放一个投票副本，每个地域都有入口和路由缓存。支付、事务恢复和发布工作者可以在多个地域运行，它们先读取持久责任，再取得当前执行资格。

\CTwentyOneFigure{ch21-01}{业务组和控制组分别跨域复制，外部平台与履约域沿自己的路径保存效果}{fig:ch21-regions}

\subsection{从放置表读出完整业务路径}

地域 A、B、C 表示设计中的独立故障域。放置表还要展开每个地域内部的机架、电力、存储和网络依赖。把 G1-A、G2-A 与控制副本装在同一台机器上，可能适合小型演练，却会使一个主机故障同时影响三个组。三个地域如果共用一个身份认证入口、一套 DNS 变更权限或一个备份账号，也会引入共同的恢复依赖。计算进程分散以后，密钥、持久卷、路由和管理权限沿哪条路径获得，同样决定灾难中的可用性。

设 Alice 的账户 1 属于 b1，组合商品包含 b2 中的 SKU-A 和 b3 中的 SKU-B。G1 的主副本在 A，G2 在 B，G3 在 C。入口 A 先把完整身份
\[
S=(1,\ \texttt{purchase.create.external.v2},\ \texttt{E1},\ \texttt{K1})
\]
交给 G1。G1 保存 $A_J$，固定 \code{P:b1:1}、\code{C:b1:1}、\code{T:b1:1} 和原支付输入。协调者再联系 G2、G3，提交联合库存的持久决定。恢复者核验 commit，推进全部必需参与者的 finish 并保存完成事实，才把原 C/T 的 \code{reservation\_ready} 置为 true。付款工作者取得这份资格以后，向原平台范围和 \code{pay-C:b1:1} 发送 3000 分的请求。平台结果返回 G1，G1 在同一次提交中形成 $R_J$、订单和 EV。消息和履约继续各自的后半段链路。

这份购买的关键路径经过三个组及平台。G1 保有多数时可以提交 b1 的受理和协调事实，SKU-A 的准备投票则由 G2 的多数保存。G2 若失去多数，SKU-A 的新预留暂停，这份购买在已保存的 blocked 付款责任处等待。G2 恢复以后，原 TX 接着推进。G1 中其他购买若只引用健康组的资源，可以继续完成各自的预留与付款。

\CTwentyOneFigure{ch21-18}{组合购买的地域路径按受理、库存决定、平台收款和事件传播依次跨过保存域}{fig:ch21-regional-path}

副本数还要与外部服务的位置一起读。支付平台可能由供应商在多个地域提供入口，也可能要求固定商户区域与出口白名单。恢复工作者 B 能否查询原收款，取决于商户身份、原对象可见范围、网络出口和凭据。把流量切到另一个商户账号会改变平台范围；新账号下相同的键可能表示另一份操作。地域切换保留原 provider、商户、测试或真实环境、动作版本和平台键，配置系统只改变到达该范围的合法网络路径。

履约也有两层位置。消费事务库保存 I/B/F 与 dispatch 意图，发货平台保存包裹或发运事实。消费者跨域恢复以后，从原 EV 和业务守卫找到 F501，再沿原 dispatch 身份查询发货结果。仓库实际位于 A，可能使发货暂停数小时，消费数据库在 B 恢复则可以先向用户报告履约已被接受。可查询、可受理和实际可交付分别对应链路中的不同终点。

\subsection{跨域读写与多数确认的成本}

入口就近部署缩短调用方到入口的距离，受理写入由所属组主副本的位置决定。三副本组的主副本收到命令后，至少等到自己与一个远端副本形成持久多数。客户端响应还要等应用完成并取到相应结果。读原结果若要求确认当前权威，可以走主副本及其有效读屏障；允许陈旧的进度投影可以就近读取，并携带观察版本与时间。前者回答原操作的确定结果，后者帮助展示最近所见进度。

设入口 A 到 G1-A 的往返为 2 毫秒，A--B 为 24 毫秒，A--C 为 68 毫秒，B--C 为 52 毫秒，各副本完成本次日志持久写入约 2 毫秒，组内应用约 1 毫秒。所选 HashiCorp Raft 实现先完成主副本的 \code{StoreLogs}，再通知复制工作者。主副本在 A、B 健康时，按这条串行路径估算，本次受理约为 $2+2+24+2+1=31$ 毫秒，分别对应入口往返、主副本保存、A--B 复制往返、B 保存和组内应用。排队、系统调度和报文处理另外进入预算。

B 故障后，G1-A 可以与 C 形成多数，复制阶段便落到 $68+2$ 毫秒附近。系统继续保持相同的受理语义，延迟预算却发生变化。若主副本改到 B，A 中入口访问 B 要经过跨域链路，B 再选择能形成多数的复制路径。主副本放置是在访问分布、幸存多数和恢复时延之间分配成本。对分区内大部分账户都在 A 的业务，稳定地把主副本放在 A 可能合适；将所有组主副本都集中到 A，则会让一个地域切换同时触发多组发现、缓存刷新和恢复突发。

组合库存还增加跨组通信。G1 向 G2、G3 并行 prepare，等待最慢的有效投票，随后保存持久决定并传播给参与者。在单组受理约 31 毫秒以后，关键路径还要依次等待库存组的持久投票、协调组的决定提交、全部必需参与者完成以及协调组保存 ready，再进入平台调用。即使每次持久边界只花几十毫秒，平台发送前也会累计可见等待。部署者沿实际依赖图并行发送库存请求，按最慢必需分支安排预算，再选择领导者位置。库存协调和付款接口即使位于同一入口进程，这些远端确认也会依次进入该进程的等待时间。

只读请求的分布也影响容量。热门订单的原结果已经不可覆盖，可以在带完整身份和结果版本的受控缓存中复用；首次观察未知结果则回到权威。落后的跟随者可能在新主已经保存 S 以后，继续返回“没有此 S”。入口随后向当前权威提交原 S 与原输入，唯一绑定命中已经提交的购买，读取与写入之间的竞态便回到同一份事实。

\subsection{链路隔离与少数侧行为}

图~\ref{fig:ch21-links} 把复制、入口、控制和平台四类通信分开。切断 A 到 B、C 的 Raft TCP，A 中旧主失去多数，B、C 可以形成新多数。若入口 HTTP 依然指向 A，用户首先见到的可能是非主提示、重定向或提交超时。入口刷新当前领导者并沿原 $S$ 重试，才能使用已经恢复的多数侧。诊断 HTTP 保持可达，有助于区分旧主的拒绝行为与端口彻底不可达。

\CTwentyOneFigure{ch21-02}{链路隔离分别作用于复制、入口、控制与平台，恢复依据沿具体通信路径判断}{fig:ch21-links}

一次超时可能覆盖两条不同历史。其一，命令在隔离前已经提交，响应在返回时丢失，新主保存了原绑定。其二，命令只到达旧主，未形成多数，新主尚未受理。两条历史给入口的观察都是超时。调用方携带同一 $S$ 与原输入访问新主，权威根据保存的提交历史处理：第一条历史读取原绑定并回放原操作，第二条历史建立首次受理。

少数侧可以继续提供明确标注的诊断和陈旧观察，但新的资源决定受本组提交能力约束。工作者若只凭内存里曾经获得的资格继续领取新 T，就会把复制隔离扩散到外部效果。实际发送前的领取或资格检查落在当前权威中，代次随责任保存。旧工作者携带过时代次回填时，权威拒绝相应资源写入。资格判断和新任期共同约束本地可接受的动作，已经发出的 HTTP 则沿下一条链路处理。

控制面暂时不可达时，已缓存的稳定归属有实际价值。入口持有 b1 属于 G1、归属 epoch 为 1 的配置，把请求送到 G1。G1 根据自己保存的激活状态核验归属。配置未变且数据组健康时，这条路径可以继续。需要首次发现、启动迁移或取得新配置资格的请求等待控制面恢复。运维暂停新的归属改变，分别观察旧配置请求在 G1 的接受情况与入口刷新失败情况，确定哪些分区继续服务、哪些请求正在等待配置。

控制面恢复以后，入口从完整快照或版本化观察重建路由，再继续接收更新。漏掉一次通知时，重新读取当前归属比只等待下一次通知更可靠。迁移完成以后，旧 G1 收到旧缓存请求，依其冻结或已交接状态拒绝旧位置写入，入口据此刷新。控制配置告诉入口去哪里，实际资源端决定这次写入是否具备资格。

\subsection{地域切换怎样处理外部在途效果}

再取一份独立初态：Alice 在平台有 10000 分，购买金额为 3000 分，G1 已保存 P1/C1/T1 与原 \code{scope/pay-C1}。旧工作者 A 在代次 7 发出付款，请求滞留在网络中。A 随后与权威多数隔离，新工作者 B 从多数侧把 T1 接管到代次 8。B 看见的本地状态是付款未知，于是按原键查询。

\CTwentyOneFigure{ch21-03}{旧付款可以在切换后到达平台，新旧观察通过原平台身份汇合到 X1}{fig:ch21-inflight}

若 B 的查询先到达平台并得到未发现，滞留的旧请求可能在查询以后到达并提交，已经丢弃的旧请求则会使平台继续保持空记录。B 此时得到的查询响应相同，后续平台状态却可能不同。若平台在原范围内按相同键和相同输入仲裁全部请求，而且原键处于保证期内，B 可以用原键恢复，平台将两次发送汇合为 X1，余额变为 7000 分。若接口只有时点查询、缺少覆盖在途请求的去重，B 就保存付款未知，收集后续回调、收款明细与人工对账结果，直到取得能关联原 C 的确定证据。

平台回调、查询响应和人工对账可能先后到达。每一份观察都携带原 C 的关联、币种、金额、账户和平台对象编号。当前有资格的恢复者验证完整输入与可信来源，再提交观察，形成原结果。A 的旧代次写入若被拒绝，X1 继续是平台已经发生的事实，B 可以重新取得并保存同一结果。运维记录被拒绝观察中的对象关联，后续审计可据此追踪。

这里需要同时保留两个时间窗口。一个是本地仍可能有迟到执行者或在途请求的时间，另一个是平台原键或可查询对象的证据保留窗口。原键只保留 24 小时的接口与长期保存收款对象的接口，适合采用的恢复动作不同。若本地在 23 小时以后才恢复，系统应优先查询并把确定对象关联持久化，发送策略由具体接口的保留保证决定。关闭某地域的出口可以减少新的请求，交换机、代理或平台侧已经接收的数据继续由原 C 承担。

\subsection{库存事务和迁移的地域依赖}

设 TX1 由 b1 协调，b2 为 SKU-A 预留 1 件，b3 为 SKU-B 预留 1 件，两份 yes 票均已持久保存，G1 此时失去多数。参与者保留 prepared 占用，C/T 保持 blocked。两份 yes 表示各库存组已准备，协调者接着要把 commit 或 abort 决定提交到自己的多数。付款门还要等全部必需参与者应用原 commit，并在协调组保存完成与 ready。协调组恢复以前，参与者超时会触发决定查询与告警，准备占用继续列入容量。

\CTwentyOneFigure{ch21-04}{协调者与参与者沿原 TX 恢复持久决定，付款资格和迁移随责任状态推进}{fig:ch21-tx}

若 commit 在 G1 故障前已经提交，新主读取到原决定，把同一决定传播给 b2、b3，参与者各自完成原子状态迁移。如果原决定尚未形成，协调恢复者按原事务规则处理保存的参与者、票据与业务条件，再提交唯一决定。全部必需参与者完成原 commit 后，协调组保存完成事实并开启原 C/T 的付款资格；abort 则关闭这份购买的付款路径。参与者与协调责任结束也使迁移准入能够通过。恢复接口把这几个状态分别展示，操作人员可以看到“已有决定、一个参与者未应用”的具体位置。

17.4 节的首版迁移等待 prepared 和协调责任结束，再冻结源分区的完整可变状态。这样可以把事务恢复保留在原部署结构中，交接时复制一个闭合图像。代价是一个持续不可达的协调组可能阻塞参与者迁移。若产品要支持带 prepared 状态迁移，就要连同稳定 TX、参与者子身份、决定查询地址、执行代次和相关日志一起转移，并给旧位置提供明确的转向规则。仅把库存数量复制过去，会让新位置缺少解释占用的持久证据。

归属 epoch、Raft 任期、工作代次和 TX 决定分别描述不同层次。b1 迁到新组以后，图像带入 P1/C1、原平台键、X1 的关联、付款未知状态及编号分配位置。新组从该位置继续分配编号，恢复器先按原键查询未知付款，再保存确定结果。若漏掉未知状态，恢复器可能把已发送责任当作首次发送；若漏掉分配位置，新购买可能误用旧外部编号。普通成员切换接续完整历史，整个集群从旧备份重建则要处理 21.3 中的历史回退与外部证据恢复。

\subsection{恢复目标落实到能力和保存点}

恢复点目标 RPO 描述允许恢复到多旧的事实，恢复时间目标 RTO 描述从故障到某项服务能力恢复的预算。一次多数副本切换若保住全部已提交日志，可以对该组已确认事实设零丢失目标；从异步灾备副本或较旧备份重建，则要按实际保存边界说明落后范围。平台可能比本地恢复点更新，两个保存域的事实必须重新关联。

恢复预算分别给原结果查询、原键查询回填、新购买受理、消息传播和实际履约设定计时终点。以下数字是一份教学预算：发现与隔离 20 秒，组内选主及路由更新 15 秒，平台出口和凭据验证 25 秒，关键未决核对 60 秒，逐步开放新流量 30 秒。若这些步骤顺序执行，新受理目标为 150 秒。平台未决积压若再花 10 分钟排空，责任清空的终点便晚于新受理开放，值守界面分别显示这两个进度。

预算还要写出并行关系。三个组可同时恢复，多组依赖的组合购买等待其中最晚恢复的必需组。下载备份与建立新机器可并行，外部核对则可能等待原商户密钥到位。调度这张依赖图，能找出最值得预置的恢复资源，例如在 B 预先验证可用的原商户出口，或者让消费者数据库具备独立的恢复目标。

真实多机部署为每个节点固定 ID、复制地址、HTTP 地址和持久目录。首次创建时，每组只有指定引导者初始化，重启已有组从原目录恢复。各副本使用自己的持久卷，备份保存快照及能衔接的日志。复制、业务和诊断接口有分别授权的访问路径。部署先建立复制组并验证成员和持久性，再写入初始归属，随后启动入口和只读观察，最后启用各类外部工作者。

十二个节点进程的放置决定故障怎样扩散。同机进程各有 TCP 连接与持久目录，主机退出会同时带走它们的计算能力；分到三台机器以后，单机退出影响各组在该机上的副本，其余机器可以重新形成多数。跨物理地域部署还把路由、区域设施和管理依赖放入切换路径。恢复计时从具体故障事件开始，连接记录标明隔离范围，图像与日志位置标明保存点，查询、受理和责任排空各自记录能力恢复时刻。
